\documentclass[11pt]{article}
\usepackage{amsmath,amssymb,amsthm,booktabs}
\usepackage[margin=1in]{geometry}
\usepackage{hyperref}

\newtheorem{lemma}{Lemma}
\newtheorem{proposition}{Proposition}
\newtheorem{theorem}{Theorem}
\newtheorem{remark}{Remark}
\newtheorem{conjecture}{Conjecture}

\title{The Prime-Indicator Matrix Does Not Obey the Semicircle Law:\\
Disproof of TLMC Conjecture 00000000142}
\author{TLMC submission (machine-verified: Lean 4, zero axioms)}
\date{September 2026}

\begin{document}
\maketitle

\begin{abstract}
TLMC Conjecture~00000000142 asserts that the empirical spectral distribution (ESD) of
the prime-indicator matrix $M_n=(\mathbf 1_{i+j\ \mathrm{prime}})_{1\le i,j\le n}$
tends to the semicircle law, with a CLT for the linear spectral statistics.
We refute the conjecture. The matrix is symmetric with $0/1$ entries, whence two exact
trace identities: $\operatorname{tr}M_n=1$ for all $n$, and
$\operatorname{tr}M_n^2=N(n):=\#\{(i,j)\in[n]^2:\ i+j \text{ prime}\}$.
Since $N(n)\sim n^2/\log(2n)$ (prime number theorem with partial summation), the ESD
second moment is $m_2(n)=N(n)/n\sim n/\log(2n)\to\infty$, while the semicircle law has
second moment $1$; the conjecture's own LSS--CLT clause presupposes $m_2(n)\to 1$.
Under the normalized reading, $\|M_n/n\|_{\mathrm{op}}\le\sqrt{N(n)}/n\to0$, so the
limit is $\delta_0$, not the semicircle. Numerically the raw spectrum spreads over the
scale $\pm\sqrt{n/\log n}$ (RMS $7.2$ at $n=300$, $16.3$ at $n=2000$), so no tight
limit exists. All finite-$n$ identities used are machine-checked in bare Lean~4 with
zero axioms.
\end{abstract}

\section{The conjecture}

\begin{conjecture}[TLMC 00000000142]
The empirical spectral distribution of the matrix
$M_n=(\mathbf 1_{i+j\ \mathrm{prime}})_{1\le i,j\le n}$ tends to the semicircle law,
and the linear spectral statistics satisfy a CLT.
\end{conjecture}

Because every prime exceeding $2$ is odd, $M_{ij}=1$ forces $i+j$ odd (except the
single entry $(1,1)$): $M_n$ is a bipartite (block off-diagonal) matrix up to the rank-one
perturbation $e_{11}e_{11}^{\!\top}$, hence by Weyl's inequality its spectrum is
symmetric about $0$ up to a $1/n$ perturbation; all eigenvalues are real.

\section{Exact trace identities}

\begin{lemma}[Trace one]\label{lem:trace}
For every $n\ge 1$, $\operatorname{tr}M_n=1$.
\end{lemma}
\begin{proof}
$\operatorname{tr}M_n=\#\{i\le n:\ 2i \text{ prime}\}$. If $i\ge2$ then $2i$ is even and
at least $4$, hence composite; and $2\cdot1=2$ is prime. So the count is $1$.
\end{proof}

\begin{lemma}[Second moment]\label{lem:tr2}
For every $n\ge1$, $\operatorname{tr}M_n^2=N(n):=\#\{(i,j)\in[n]^2:\ i+j \text{ prime}\}$.
\end{lemma}
\begin{proof}
$M$ is symmetric with entries in $\{0,1\}$, so
$\operatorname{tr}M^2=\sum_{i,j}M_{ij}M_{ji}=\sum_{i,j}M_{ij}^2=\sum_{i,j}M_{ij}=N(n)$.
\end{proof}

\begin{proposition}[Asymptotics of $N(n)$]\label{prop:asymp}
$N(n)\sim n^2/\log(2n)$ as $n\to\infty$. Consequently the ESD second moment
$m_2(n)=\frac1n\operatorname{tr}M_n^2=N(n)/n\sim n/\log(2n)\to\infty$.
\end{proposition}
\begin{proof}[Proof sketch]
$N(n)=\sum_{p\le2n}g_n(p)$ with $g_n(p)=\min(p-1,\,2n+1-p)$. Since
$g_n(p)=n\,\varphi(p/n)$ with $\varphi(t)=\min(t,2-t)$ a fixed continuous function on
$[0,2]$ and $\int_0^2\varphi(t)\,dt=1$, partial summation with the prime number theorem
$\pi(x)\sim x/\log x$ gives $N(n)\sim n^2/\log(2n)$. (Convergence is logarithmically
slow; numerically $N(n)\log(2n)/n^2\approx1.09$--$1.11$ for $300\le n\le8000$, drifting
toward $1$.)
\end{proof}

\section{Refutation}

\begin{theorem}[Raw reading]\label{thm:raw}
If the ESD of $M_n$ tends to the semicircle law and the linear spectral statistics
satisfy a CLT in the standard (polynomial test function) sense, then in particular
$\frac1n\operatorname{tr}M_n^2\to\int x^2\,d\sigma_{\mathrm{sc}}(x)=1$.
But $\frac1n\operatorname{tr}M_n^2=N(n)/n\to\infty$ by
Lemmas~\ref{lem:tr2} and Proposition~\ref{prop:asymp}. Contradiction.
\end{theorem}

\begin{theorem}[Normalized reading]\label{thm:norm}
$\|M_n/n\|_{\mathrm{op}}\le\|M_n\|_F/n=\sqrt{N(n)}/n\le\sqrt{\pi(2n)/n}\to0$ by the
prime number theorem. Hence every eigenvalue of $M_n/n$ lies in
$[-\sqrt{2/\log(2n)},\sqrt{2/\log(2n)}]$ eventually and the ESD of $M_n/n$ converges to
$\delta_0$, not to the semicircle law.
\end{theorem}

\begin{remark}[Relation to the recorded attack]
The recorded verdict recomputed $\operatorname{tr}M_n^2/n^2$: $0.172$, $0.156$, $0.145$
at $n=300,600,1000$ and concluded the limit is $\delta_0$. Our independent recomputation
confirms those numbers exactly ($0.1715$, $0.1562$, $0.1452$). One normalization is
corrected here: $\operatorname{tr}M^2/n^2$ is the squared Frobenius density, not the ESD
second moment, which is $\operatorname{tr}M^2/n$ and diverges; the honest $\delta_0$
statement is Theorem~\ref{thm:norm}, for the eigenvalue-normalized matrix $M_n/n$. Under
either normalization the conjecture is false --- in the raw reading it fails more
strongly than recorded, since no tight limit exists at all.
\end{remark}

\begin{remark}[What carries the raw-reading refutation]
Weak limits are not determined by moments without uniform integrability; the divergence
of $m_2(n)$ alone would leave the logical loophole of mass escaping to infinity. Here
the conjecture's second clause closes it: a CLT for linear spectral statistics with
polynomial $f$ entails the centered moment convergence used in
Theorem~\ref{thm:raw}, and that fails outright. Independently, the numerics of
Section~\ref{sec:num} show the bulk itself (not a vanishing fraction) at scale
$\sqrt{n/\log n}$, so no distributional limit of the ESD of $M_n$ exists, tight or not.
\end{remark}

\section{Numerical verification}\label{sec:num}

Exact integer counts $N(n)$ (sieve; also fixed by \texttt{decide} in Lean for
$n=8,64,128$: $23$, $941$, $3239$):

\begin{center}
\begin{tabular}{rrrrrr}
\toprule
$n$ & $N(n)$ & $N(n)/n^2$ & $1/\log(2n)$ & $N(n)/n$ (ESD $m_2$) & $n/\log(2n)$\\
\midrule
300 & 15439 & 0.1715 & 0.1563 & 51.46 & 46.90\\
600 & 56237 & 0.1562 & 0.1410 & 93.73 & 84.63\\
1000 & 145171 & 0.1452 & 0.1316 & 145.17 & 131.56\\
2000 & 528537 & 0.1321 & 0.1206 & 264.27 & 241.14\\
4000 & 1944355 & 0.1215 & 0.1113 & 486.09 & 445.08\\
8000 & 7187475 & 0.1123 & 0.1033 & 898.43 & 826.42\\
\bottomrule
\end{tabular}
\end{center}

Spectral statistics (symmetric eigensolver; \texttt{reproduce.py --eigen}):

\begin{center}
\begin{tabular}{rrrrrrr}
\toprule
$n$ & $\lambda_{\max}$ & $\lambda_{\min}$ & RMS & $\tfrac1n\#\{|\lambda|\le1\}$ &
$\tfrac1n\#\{\lambda<0\}$ & $\tfrac1n\#\{|\lambda|>3\}$\\
\midrule
300 & 51.72 & $-51.71$ & 7.17 & 0.1200 & 0.5000 & 0.5767\\
600 & 94.03 & $-94.03$ & 9.68 & 0.0667 & 0.5000 & 0.7400\\
1000 & 145.62 & $-145.61$ & 12.05 & 0.0520 & 0.5000 & 0.7860\\
2000 & 264.96 & $-264.95$ & 16.26 & 0.0320 & 0.5000 & 0.8690\\
\bottomrule
\end{tabular}
\end{center}

Semicircle reference values: second moment $1$;
$\sigma_{\mathrm{sc}}([-1,1])=\tfrac13+\tfrac{\sqrt3}{2\pi}\approx0.609$;
$\sigma_{\mathrm{sc}}(\mathbb R\setminus[-3,3])=0$. The negative fraction is exactly
$1/2$ at all $n$ (bipartite symmetry), the only statistic that agrees. Note
$\lambda_{\max}\approx N(n)/n$: the Rayleigh lower bound
$\lambda_{\max}\ge\mathbf 1^{\!\top}M_n\mathbf 1/n=N(n)/n$ is nearly tight, so the
outlier pair $\pm\lambda_1$ alone accounts for the bulk of $\operatorname{tr}M_n^2$.

\section{Boundary: proven, heuristic, open}

\textbf{Proven and machine-checked (Lean 4, zero axioms):} Lemma~\ref{lem:trace} and
Lemma~\ref{lem:tr2} for all $n$; the exact values $N(8)=23$, $N(64)=941$,
$N(128)=3239$; $N(n)>n$ at $n=64,128$ (ESD second moment above the semicircle's total
second moment); the growth witness $N(128)/128>N(64)/64$; $\operatorname{tr}M_{1000}=1$
by kernel computation.

\textbf{Proven at PNT level (cited, not formalized):} $N(n)\sim n^2/\log(2n)$;
$m_2(n)\to\infty$; $\|M_n/n\|_{\mathrm{op}}\to0$ and $\mathrm{ESD}(M_n/n)\to\delta_0$.

\textbf{Heuristic (stated as such):} the bulk of the spectrum of $M_n$ spreads over a
symmetrized quarter-circle / Marchenko--Pastur type shape at scale
$\approx\sqrt{2n/\log(2n)}$ (the bipartite block $A$ has density $\approx2/\log(2n)$ and
the rank-one mean contribution escapes as the outlier pair $\pm n/\log(2n)$).

\textbf{Open:} a rigorous identification of the rescaled bulk limit; a proof that a
\emph{positive fraction} of eigenvalues escapes any fixed compact (Frobenius counting
gives only $\gtrsim c\log n$ eigenvalues outside $[-K,K]$); linear spectral statistics
of any rescaled model.

\section{Reproduction}

\begin{verbatim}
python3 reproduce.py            # exact counts, trace identities, verdict comparison
python3 reproduce.py --eigen    # + eigenvalue statistics (numpy)
cd lean4 && lake build && lake env lean Check.lean
\end{verbatim}

\section{Machine verification}

Bare Lean~4 (toolchain \texttt{leanprover/lean4:v4.33.1}, no Mathlib,
\texttt{lakefile.toml} with \texttt{defaultTargets = ["Main"]}). \texttt{Check.lean}
prints \texttt{\#print axioms} for all 18 theorems; every line reads
\texttt{does not depend on any axioms} (no \texttt{sorryAx}, no \texttt{propext}, no
\texttt{Quot.sound}, no \texttt{Classical.choice}). The proofs use only
\texttt{rfl}, \texttt{decide} (kernel computation), explicit
\texttt{Eq.trans}/\texttt{congrArg}/\texttt{Eq.subst} term proofs and structural
induction; \texttt{omega}, \texttt{simp} and \texttt{native\_decide} are deliberately
avoided because they would introduce axioms.

\begin{thebibliography}{9}
\bibitem{wigner} E. Wigner, Characteristic vectors of bordered matrices with infinite
dimensions, \emph{Ann. of Math.} \textbf{62} (1955), 548--564.
\bibitem{mp} V. A. Mar\v{c}enko and L. A. Pastur, Distribution of eigenvalues for some
sets of random matrices, \emph{Math. USSR-Sb.} \textbf{1} (1967), 457--483.
\bibitem{pnt} J. Hadamard (1896); Ch. de la Vall\'ee Poussin (1896), on the prime
number theorem $\pi(x)\sim x/\log x$.
\bibitem{bai} Z. Bai and J. W. Silverstein, \emph{Spectral Analysis of Large
Dimensional Random Matrices}, Springer, 2010.
\bibitem{taovu} T. Tao and V. Vu, Random matrices: universality of local spectral
statistics of non-Hermitian matrices, \emph{Ann. Probab.} \textbf{43} (2015), 782--874.
\end{thebibliography}

\end{document}
